\chapter{Electric potential, conductors and capacitance}\label{ch:pot}

The electric field is a conservative one, so chapter~4's energy grammar
applies: define potential energy, divide by charge, and get $V$ --- the
scalar field whose gradient is $\mathbf E$. Capacitance is the
engineering payout.

\section{Potential energy and the sign discipline}\label{sec:pot-u}

For two charges $U=kq_1q_2/r$ --- positive for likes (work in to
assemble), negative for unlikes (bound). Moving $q$ through
$\Delta V$ costs $W=q\Delta V$ with the same sign rules as work
everywhere.

\section{$V$ and equipotentials}\label{sec:pot-v}

\begin{equation}
  V(\mathbf r)=k\sum_i\frac{q_i}{r_i},\qquad
  \mathbf E=-\nabla V,\label{eq:v}
\end{equation}
a scalar field --- no vector sums. Equipotentials are level sets and
cross field lines at right angles (figure~\ref{fig:equip}); conductors
are single equipotentials, which is why the surface field is normal to
the metal.

\begin{figure}[htbp]\centering
\begin{tikzpicture}[scale=0.9,font=\scriptsize]
\foreach \r in {0.6,1.1,1.7,2.4}{\draw[accent] (0,0) circle (\r);}
\foreach \ang in {0,45,...,315}{\draw[->,thick,amber] (0,0) -- ({\ang}:2.7);}
\fill[amber] (0,0) circle (0.07);
\node at (0,-3.0) {equipotentials (circles) cross field lines (rays) at right angles};
\end{tikzpicture}

\caption{Equipotential circles and radial field lines for a point
charge --- always orthogonal.}
\label{fig:equip}
\end{figure}

\begin{example}[Assembling a configuration]\label{ex:assemble}
Place $+q$ at $(-1,0)$, $+q$ at $(1,0)$, $-q$ at $(0,1)$. Find the
assembly energy.
\par\smallskip
\textbf{Strategy.} Bring charges one at a time; each pays $q_iV$ of
those already placed.
\begin{solution}
First free; second pays $kq^{2}/2$; third sees both:
$-kq^{2}/\sqrt2-kq^{2}/\sqrt2=-\sqrt2\,kq^{2}$.
Total $U=kq^{2}(1/2-\sqrt2)=-0.914\,kq^{2}$ --- bound overall.
The code below verifies both the sequential and pairwise sums.
\checkmark
\end{solution}
\end{example}

\section{$\mathbf E=-\nabla V$: the gradient means steepness}\label{sec:pot-grad}

In one dimension $E_x=-dV/dx$: where $V$ slopes, fields point downhill
(figure~\ref{fig:vx}); where $V$ is flat, the field dies. Voltage is
the map; the field is its slope.

\begin{figure}[htbp]\centering
\begin{tikzpicture}
\begin{axis}[width=8cm,height=4.6cm,xlabel={$x$},ylabel={},
xmin=-3,xmax=3,ymin=-2,ymax=2,font=\small,legend pos=north east]
\addplot[very thick,accent,domain=-2.9:2.9,samples=80] {1/(abs(x)+0.3)};
\addlegendentry{$V$}
\addplot[very thick,amber,domain=-2.9:2.9,samples=80] {-x/(abs(x)+0.3)^2};
\addlegendentry{$E=-dV/dx$}
\addplot[dashed] coordinates {(-3,0) (3,0)};
\node[font=\scriptsize] at (axis cs:1.5,-1.2) {steepest $V$ = biggest $E$};
\end{axis}
\end{tikzpicture}

\caption{$V(x)$ and $E=-dV/dx$ side by side: field peaks where potential
slopes hardest.}
\label{fig:vx}
\end{figure}

\begin{example}[$V$ on the axis of a charged disk]\label{ex:disk}
A disk of radius $R$ carries $\sigma$ uniformly. Find $V(x)$ on axis
and check $E=-dV/dx$ reproduces the ring result.
\par\smallskip
\textbf{Strategy.} Integrate rings --- $V$ is scalar so there's no
component bookkeeping.
\begin{solution}
Ring of radius $r$, width $dr$: $dV=k\sigma2\pi r\,dr/\sqrt{x^{2}+r^{2}}$.
\[V=\frac{\sigma}{2\varepsilon_0}\int_0^{R}\frac{r\,dr}{\sqrt{x^{2}+r^{2}}}
=\frac{\sigma}{2\varepsilon_0}\left(\sqrt{x^{2}+R^{2}}-|x|\right).\]
Differentiating gives
$E=\frac{\sigma}{2\varepsilon_0}\left(1-\frac{x}{\sqrt{x^{2}+R^{2}}}\right)$
--- the disk's known result, and at $R\to\infty$ it returns the sheet's
$\sigma/2\varepsilon_0$. \checkmark
\end{solution}
\end{example}

\section{Capacitance: geometry's charge budget}\label{sec:pot-cap}

For any two conductors $C=Q/V$ --- a pure geometry number. Parallel
plates ($E=\sigma/\varepsilon_0$ over gap $d$) give
$C=\varepsilon_0A/d$; the coaxial and spherical cases
(figure~\ref{fig:caps}) follow the same $E\to V\to C$ chain.

\begin{figure}[htbp]\centering
\begin{tikzpicture}[scale=0.85,font=\scriptsize]
\draw[fill=accent!12,very thick] (0,0) circle (0.5);
\draw[very thick] (0,0) circle (1.1);
\node at (0,-1.5) {coaxial: $C=2\pi\varepsilon_0L/\ln(b/a)$};
\node[align=left] at (0,-1.5) {};
\draw[fill=accent!12,very thick] (4,0) circle (0.5);
\draw[very thick] (4,0) circle (1.1);
\node[align=left] at (4,-1.5) {spherical:\\$C=4\pi\varepsilon_0ab/(b-a)$};
\end{tikzpicture}

\caption{Coaxial and spherical capacitors: same three-step derivation,
different symmetry.}
\label{fig:caps}
\end{figure}

\begin{figure}[htbp]\centering
\begin{tikzpicture}[scale=0.95,font=\scriptsize]
\draw[very thick] (-1.6,1.0) -- (1.6,1.0) node[right]{$+$};
\draw[very thick] (-1.6,-1.0) -- (1.6,-1.0) node[right]{$-$};
\foreach \x in {-1.2,-0.6,0,0.6,1.2}{\draw[->,accent] (\x,0.9) -- (\x,-0.9);}
\draw[accent,dashed] (-1.8,0.9) .. controls (-2.3,0) .. (-1.8,-0.9);
\draw[accent,dashed] (1.8,0.9) .. controls (2.3,0) .. (1.8,-0.9);
\node at (0,-1.6) {uniform inside, fringing at edges ($d\ll$ plate)};
\end{tikzpicture}

\caption{The parallel-plate capacitor: uniform field inside, fringing at
the edges.}
\label{fig:plate}
\end{figure}

\begin{example}[Spherical capacitor end-to-end]\label{ex:sphcap}
Find the capacitance of concentric spheres $a<b$.
\par\smallskip
\textbf{Strategy.} $E$ by Gauss, $V$ by integrating it, $C=Q/V$.
\begin{solution}
$E=kQ/r^{2}$ between, so
$V=\int_a^{b}E\,dr=kQ(1/a-1/b)=kQ(b-a)/ab$ and
$C=4\pi\varepsilon_0ab/(b-a)$. Limit check: $b\to\infty$ leaves an
isolated sphere at $4\pi\varepsilon_0a$ --- the van de Graaff dome's
capacity. \checkmark
\end{solution}
\end{example}

\section{Dielectrics and stored energy}\label{sec:pot-diel}

A dielectric polarizes: bound surface charge opposes the free charge's
field (figure~\ref{fig:diel}), cutting $E$ by $\kappa$ and raising $C$
by the same factor. The energy stored is
\begin{equation}
  U=\tfrac12CV^{2}=\tfrac12QV=\frac{Q^{2}}{2C}
  =\frac{\varepsilon_0}{2}\int E^{2}\,dV,\label{eq:u}
\end{equation}
which says the field itself carries energy --- the seed of the
electromagnetic-wave chapters.

\begin{figure}[htbp]\centering
\begin{tikzpicture}[scale=0.9,font=\scriptsize]
\draw[very thick] (-1.6,1.1) -- (1.6,1.1) node[right]{$+$};
\draw[very thick] (-1.6,-1.1) -- (1.6,-1.1) node[right]{$-$};
\draw[fill=amber!15,very thick] (-1.4,-0.7) rectangle (1.4,0.7);
\foreach \x/\s in {-1.0/+,-0.4/+,0.2/+,0.8/+}{\node at (\x,0.15) {$\mp$};}
\foreach \x/\s in {-1.0/-,-0.4/-,0.2/-,0.8/-}{\node at (\x,-0.15) {$\pm$};}
\node[align=left] at (0,-1.7) {dielectric: bound surface charge\\weakens $E$ by $\kappa$};
\node[align=left] at (0,-1.7) {};
\end{tikzpicture}

\caption{Dielectric polarization: bound charge weakens the interior
field by $\kappa$.}
\label{fig:diel}
\end{figure}

\begin{figure}[htbp]\centering
\begin{tikzpicture}
\begin{axis}[width=8cm,height=4.4cm,xlabel={$V$},ylabel={$U$},
xmin=0,xmax=10,ymin=0,ymax=6,font=\small,legend pos=north west]
\addplot[very thick,accent,domain=0:10] {0.06*x^2};
\addlegendentry{$U=\frac12CV^2$}
\addplot[dashed,amber] coordinates {(0,0) (10,6)};
\node[font=\scriptsize,align=left] at (axis cs:4,4.5) {charging = lifting\\charges up the hill};

\end{axis}
\end{tikzpicture}

\caption{$U=\tfrac12CV^{2}$: the charging work piles up against the
growing voltage.}
\label{fig:uenergy}
\end{figure}

\begin{example}[Series-parallel combination]\label{ex:series}
A $10$\,\textmu F and a $40$\,\textmu F capacitor are wired in series
across $100$\,V. Find each voltage; then the parallel equivalent.
\par\smallskip
\textbf{Strategy.} Series capacitors share charge, not voltage.
\begin{solution}
$C_s=(1/10+1/40)^{-1}=8$\,\textmu F, so $Q=C_sV=800$\,\textmu C on
each: $V_1=80$\,V, $V_2=20$\,V --- the small one takes the big share.
Parallel gives $50$\,\textmu F at the common voltage. Rule of thumb:
series shares $Q$ (small wins the voltage), parallel shares $V$.
\checkmark
\end{solution}
\end{example}

\begin{example}[Camera-flash energy audit]\label{ex:flash}
A $100$\,\textmu F capacitor charges to $300$\,V. Find $Q$, $U$, and
the average power if the flash dumps it in $1$\,ms.
\par\smallskip
\textbf{Strategy.} $Q=CV$; $U=\tfrac12CV^{2}$; $P=U/t$.
\begin{solution}
$Q=0.030$\,C; $U=\tfrac12(10^{-4})(9\times10^{4})=4.5$\,J;
$P=4.5/0.001=4.5$\,kW --- four kilowatts for a millisecond from a
battery that trickled seconds charging it. Capacitors are time-traders.
\checkmark
\end{solution}
\end{example}

\begin{lstlisting}[language=Python,caption={\texttt{code/ch11\_potential\_map.py} --- gradient check and assembly-energy audit.},label={lst:pot}]
dVdx = (V(charges, x+h, y) - V(charges, x-h, y))/(2*h)
\end{lstlisting}

Output of a real run:
\begin{lstlisting}[numbers=none]
Check -dV/dx ~ Ex at sample points (k=1):
  (1.2,0.3): -dV/dx=-1.2543  Ex=-1.2543  match=True
  (0.8,-0.4): -dV/dx=-1.8834  Ex=-1.8834  match=True
  (2.0,0.8): -dV/dx=-0.1671  Ex=-0.1671  match=True

Assembly work step by step (sum q_i*V_others):
  bring +1.0 to (-1,0): W=0 (first)
  bring +1.0 to (1,0): W=+0.5000
  bring -1.0 to (0,1): W=-1.4142
  total assembly energy U = -0.9142 (=sum k q_i q_j/r_ij)
  pairwise cross-check: -0.9142
\end{lstlisting}

The finite difference reproduces $\mathbf E$ to all printed digits, and
the two energy routes --- sequential assembly vs pairwise sum --- agree
to $-0.9142$ as they must.

\section{In engineering practice: every capacitor on the board}\label{sec:pot-eng}

Decoupling capacitors are \eqref{eq:u} as a local battery (fast charge
right where the chip needs it), supercaps trade area $A$ for farads via
electrochemical surfaces, and the camera flash of
example~\ref{ex:flash} is the canonical time-trade. Cable capacitance
--- the coaxial formula of \S\ref{sec:pot-cap} --- is why long signal
runs need drivers: hundreds of pF per metre must be recharged every
edge.

\section{Common mistakes}\label{sec:pot-mistakes}

\begin{tripwire}[Sign discipline dropped]
$U=kq_1q_2/r$ carries signs: unlike charges give \emph{negative}
energy. Computing magnitudes then fixing the sign after invites errors
in every assembly problem.
\end{tripwire}
\begin{tripwire}[$V$ treated as a vector]
$V$ is scalar: sum it, don't vector-add it. The gradient then gives
$\mathbf E$'s direction for free.
\end{tripwire}
\begin{tripwire}[$\Delta V$ vs $V$]
Only differences are physical for conductors (which sit at some
$V$, not ``at $V$''); the choice of zero is convention.
\end{tripwire}
\begin{tripwire}[Series/parallel swapped]
Capacitors combine like resistors inverted: series gives
$C^{-1}=\sum C_i^{-1}$ (smallest $C$ wins the voltage), parallel sums
$C$.
\end{tripwire}
\begin{tripwire}[$U=QV$ instead of $\tfrac12QV$]
The charging work is against the \emph{growing} voltage --- half the
charge rode in at lower $V$, so the total is $\tfrac12QV$, not $QV$.
\end{tripwire}

\section{Summary}\label{sec:pot-summary}

\begin{table}[ht]\centering
\begin{tabularx}{\linewidth}{@{}lY@{}}
\toprule
\tablehead\thead{Idea} & \thead{One-line version}\\
\midrule
$U$ of pair & $kq_1q_2/r$ with signs \\
$V$ & $k\sum q_i/r_i$ scalar; $\mathbf E=-\nabla V$ \\
equipotential & level set; orthogonal to $\mathbf E$; conductor = one \\
capacitor & plates $\varepsilon_0A/d$; sphere $4\pi\varepsilon_0ab/(b-a)$ \\
dielectric & $C$ multiplies by $\kappa$ \\
energy & $\tfrac12CV^{2}$ or $\frac{\varepsilon_0}{2}\int E^{2}dV$ \\
\bottomrule
\end{tabularx}
\caption{Potential and capacitance cheat sheet.}
\label{tab:potsummary}
\end{table}

\begin{wins}[Definition of done]
You can now: sum $V$ scalar-wise, differentiate it for $\mathbf E$,
audit assembly energy, derive $C$ for plate/coax/sphere, and price
stored energy.
\end{wins}

\chapternotes{%
The scalar-first treatment follows MIT 8.02SC \cite{ocw801} and AP
Physics C E\&M \cite{apced}; Purcell--Morin's potential chapter is the
reference for the $\varepsilon_0E^{2}/2$ energy-density view
\cite{hrw}. The numeric checks are the listing's real output.}

\begin{rulebox}[{Watch instead}]
\begin{enumerate}
\item \emph{Electrostatic Potential, Electric Energy, Equipotential
Surfaces} --- Walter Lewin, MIT 8.02x Lecture 4, 49:01,
\url{https://youtu.be/QpVxj3XrLgk} ---
\S\ref{sec:pot-v}--\S\ref{sec:pot-grad} at full length (accessed
2026-10-11).
\item \emph{Capacitors and capacitance} --- Khan Academy, 5:43,
\url{https://youtu.be/u-jigaMJT10} --- \S\ref{sec:pot-cap} worked
quickly (accessed 2026-10-11).
\end{enumerate}
\end{rulebox}

\section*{Further reading}
\addcontentsline{toc}{section}{Further reading}

\begin{itemize}
\item MIT OCW \emph{8.02SC}, potential and capacitance modules. At
\url{https://ocw.mit.edu}, course ``8.02SC'' (verified 2026-10-11).
\item OpenStax \emph{College Physics 2e}, chs.~19--20 --- free text:
\url{https://openstax.org} (verified 2026-10-11).
\item College Board, \emph{AP Physics C (E\&M) CED}, Units 8--9 ---
verified 2026-10-11 at \url{https://apstudents.collegeboard.org}.
\item E.~M. Purcell and D.~J. Morin, \emph{Electricity and Magnetism},
chs.~2--3 --- potential and capacitance with the energy-density
derivation.
\item R.~P. Feynman, \emph{Lectures on Physics}, vol.~II, ch.~8 ---
electrostatic energy as field energy; free at Caltech's site.
\item Khan Academy, \emph{Circuits: capacitors} --- drills matching
\S\ref{sec:pot-cap}--\S\ref{sec:pot-diel}, at
\url{https://www.khanacademy.org} (verified 2026-10-11).
\end{itemize}

\section*{Additional exercises}
\addcontentsline{toc}{section}{Additional exercises}

\begin{exercise}[easy]
Two charges $+2$\,\textmu C and $-1$\,\textmu C sit $1$\,m apart:
find $V$ at the midpoint and the work to move a third charge there.
(AP-C)
\end{exercise}
\begin{exercise}[easy]
From an equipotential map, two lines $10$\,V apart are $2$\,cm apart:
estimate $|E|$. (AP-C)
\end{exercise}
\begin{exercise}[easy]
Combine $10$\,\textmu F and $30$\,\textmu F in series and in parallel.
(AP-C)
\end{exercise}
\begin{exercise}[easy]
A $47$\,\textmu F capacitor at $12$\,V: find $Q$ and $U$. (F=ma)
\end{exercise}
\begin{exercise}[easy]
Find $V$ a distance $r$ from a point charge $-Q$ and explain the sign.
(AP-C)
\end{exercise}
\begin{exercise}[medium]
A ring carries $Q$ on radius $R$: find $V(x)$ on axis, then
differentiate for the ring's $E$ --- check against ch9. (AP-C FR-style)
\end{exercise}
\begin{exercise}[medium]
How much work to bring a fourth charge $+q$ to the centre of the
configuration of example~\ref{ex:assemble}? (F=ma)
\end{exercise}
\begin{exercise}[medium]
A parallel-plate capacitor halves its plate separation while isolated:
what happens to $C$, $V$, $U$? (AP-C)
\end{exercise}
\begin{exercise}[medium]
Derive the coaxial capacitor formula $C=2\pi\varepsilon_0L/\ln(b/a)$.
(AP-C)
\end{exercise}
\begin{exercise}[medium]
Insert a $\kappa=3$ dielectric into a charged isolated capacitor:
describe what changes and compute the new energy. (AP-C)
\end{exercise}
\begin{exercise}[medium]
Two capacitors in series share $100$\,V: find each voltage for
$10$\,\textmu F and $40$\,\textmu F. (engineering)
\end{exercise}
\begin{exercise}[hard]
Find the energy stored in a charged sphere of radius $R$ by integrating
$\varepsilon_0E^{2}/2$ over space, and match $\tfrac12Q^{2}/C$.
(USAPhO)
\end{exercise}
\begin{exercise}[hard]
A dipole sits at angle $\theta$ in a uniform field: find the torque
from $U=-\mathbf p\cdot\mathbf E$ and the oscillation frequency for
small $\theta$. (APhO-style)
\end{exercise}
\begin{exercise}[medium]
Show the field energy inside a parallel-plate capacitor equals
$\tfrac12CV^{2}$ directly from $\varepsilon_0E^{2}/2\times$volume.
(AP-C)
\end{exercise}
\begin{exercise}[medium]
An electron volt is one electron through one volt: convert $1$\,eV to
joules and give the proton's speed at that energy. (F=ma)
\end{exercise}
